\section{解各种不等式}

\subsection{解绝对值不等式}

\begin{theorem}[绝对值不等式]
    使用零点分段讨论法，将绝对值不等式分段书写。
\end{theorem}

\begin{example}
    令 $y=|x-1|+2|x+1|$，在平面直角坐标系（下图）中画出函数 $y=|x-1|+2|x+1|$ 的图象。并解不等式 $y \leq 5$
\end{example}

\v{10em}

\begin{example}
    证明：$|x+2|-|x-1| \geq-3$ ，对所有实数 $x$ 均成立，并求等号成立时 $x$ 的取值范围．
\end{example}

\vspace{10em}

\subsection{解一元二次不等式}

\begin{theorem}[一元二次不等式]
    实际上就先解一元二次方程，根据图像查看不等式的取值范围即可。
\end{theorem}

\begin{example}
    解不等式：$x^2-2x-3>0$
\end{example}

\v{6em}

\subsubsection{一元二次不等式的二次系数讨论}

\begin{theorem}[二次系数讨论]

\end{theorem}

\subsubsection{一元二次不等式的含参因式分解}

\begin{theorem}[]

\end{theorem}

\subsubsection{}

\subsection{解一元高次不等式}

\begin{theorem}[数轴穿根法]

\end{theorem}